\chapter{Related Work}
Dimensional stacking is a visualization technique that allows for the simultaneous display of multiple dimensions or categories of data within a single plot. It involves stacking multiple two-dimensional plots on top of each other, with each plot representing a different category or dimension of the data. While dimensional stacking can be a useful technique for visualizing multiple dimensions of data, there are also other techniques available for multidimensional visualization. There will be no clarity in this technique as all dimensions are plotted instead the visualization techniques like bivaraite and univariate techniques are used in our methodology.

\section{Visualization Recommendation Approaches}
Visualization recommendation approaches can be divided into rule based, machine learning and  knowledge graph approaches. Rule based visualization recommendation approaches will be constructing the rules depending on the type of the data in the dataset and with the help of experts intervention. Examples of rule based approaches include APT \cite{mackinlay1986automating}, Show Me\cite{mackinlay2007show} SAGE \cite{roth1994interactive} and Voyager2 \cite{wongsuphasawat2017voyager}. The disadvantage is that it requires a lot of human effort for constructing the rules.  Contrasting to rule based approaches, machine learning based approaches will be extracting the features from the dataset and those extracted features are given to machine learning models for predicting the visualization of a particular dataset.  Qian et al., \cite{qian2020ml} and Hu et al., \cite{hu2019vizml} are examples of machine learning based approaches. These approaches use neural networks which is a black box where we are unable to know why the particular visualization is recommended. Knowledge graph based approaches such as KG4Vis \cite{li2021kg4vis} will construct the knowledge graph from the visualization corpus and use the graph to recommend the visualization. Constructing the knowledge graph is the main hindrance in this approach which is mainly dependent on a large corpus of the data. There are several other approaches that are related to the visualization recommendation in different contexts which are discussed below.

Lee et.al., \cite{lee2021lux} approach extracts information from the data frames and recommends the visualization. Epperson et.al., \cite{epperson2022leveraging} recommend the visualization based on the history of a user’s data analysis. Ojo et.al., \cite{ojo2022visgnn} recommend personalized visualization by constructing Graph Neural Network from the large corpus of the data. Harris et.al., \cite{harris2021insight} recommend visualization by utilizing relevant insights in the data. Gao et.al., \cite{gao2021vizgrank} obtained all possible visualizations for the given dataset, constructed a graph by modeling the relations between all visualizations, and finally recommended the high score visualization. Kelleher et.al., \cite{kelleher2021introductory} provided the visualization recommendation for scientific datasets in the context of big data. 
The above discussed approaches deal with different dataset for visualization recommendation. There exists an approach Maruta et al. \cite{maruta2021intent}  which is based on two input parameters such as intent and dataset used for visualization recommendation. Intent refers to information that actually the user desires to visualize from the dataset.



Apart from visualization recommendation approaches discussed earlier there is a need to know about incremental visualization in order to solve the problem of visualization recommendation for incremental data. Incremental visualization \cite{schulz2015enhanced} is a technique used in data visualization where data is updated dynamically in real-time. There is no efficient strategy in the literature to address incremental visualization recommendation. In this paper we proposed the methodology to solve this problem without rendering an entire chart or graph for every time stamp which can save significant processing time and improves the user experience. We extended the works of Maruta et al. \cite{maruta2021intent} which recommends the visualization based on two input parameters such as the intent and dataset provided by the user for incremental data.

\section{Exploring basic NLDR techniques}
Non-Linear Dimensionality Reduction (NLDR) approaches are indispensable for visualizing high-dimensional data, especially when linear methods are inadequate for capturing intricate data
structures. Techniques such as t-SNE \cite{tsne}, UMAP \cite{umap_reference}, and Laplacian-based methods \cite{belkin2001laplacian} \cite{sun2023laplacian} play a crucial role in preserving local and global structures by projecting data into lower dimensional
spaces while retaining meaningful patterns. These methods are particularly effective in
uncovering non-linear relationships, making them invaluable for exploring and understanding
complex datasets, especially in fields like healthcare and marketing.


An approach like Sammon Mapping \cite{henderson1997sammon} minimizes the difference between pairwise distances in the high-dimensional data and their corresponding distances in the low-dimensional representation. Curvilinear Component Analysis \cite{demartines1995cca} identifies a low-dimensional data representation that maximizes the correlation between the actual and low-dimensional data points. Isomap \cite{tenenbaum2000global} improves traditional MDS \cite{hout2013multidimensional} by including geodesic distances, making it ideal for nonlinear manifolds. SNE \cite{hinton2002stochastic} maintains the pairwise similarities between data points by applying probability distributions while transforming high-dimensional data into low-dimensional data and capturing local structures.

\section{From Neighbourhood to Topologies: Embedding High-Dimensional Data}

SNE \cite{hinton2002stochastic} minimizes the Kullback-Leibler (KL) divergence between the probability distributions in low and high-dimensional spaces. t-SNE \cite{tsne} effectively maps high-dimensional data to a low-dimensional space while preserving pairwise similarities between data points, unlike its predecessor, SNE \cite{hinton2002stochastic}, which relies on asymmetric conditional probabilities. High-dimensional similarities are modeled using Gaussian distributions, whereas low-dimensional similarities leverage a student-t distribution with one degree of freedom, providing heavy tails to capture the influence of outliers. t-SNE \cite{tsne} introduces symmetric joint probabilities, resulting in a symmetric and more stable cost function. Through gradient descent optimization, t-SNE \cite{tsne} overcomes SNE’s limitations by improving complex data handling for visualization and analysis. However, the authors claim that it preserves global structure to some extent and is computationally intensive.


McInnes et al. \cite{umap_reference} discovered UMAP, the NLDR technique, which creates a high-dimensional graph, and then a similar graph in lower dimensions is determined. UMAP relies on strong theoretical principles of riemannian manifolds \cite{lin2008riemannian} and preservation of topological structure. UMAP operates through two key phases: (1) the construction of a high-dimensional graph and (2) the optimization of low-dimensional embedding. 


In the first phase, a graph is constructed by identifying the k-nearest neighbours of each data point $x_i$ in high-dimensional space using pairwise distances. The high-dimensional probabilities are calculated to represent the fuzzy membership strengths between data points. A weighted k-nearest neighbour graph is then constructed. To symmetrize the adjacency matrix, UMAP applies a transformation, combining the matrix with its transpose.


In the second phase, the low-dimensional space is initialized randomly, and its probabilities are computed for each pair of points based on their distances, with hyperparameters \textit{`a'} (scaling factor) and \textit{`b'} (decay rate). The final low-dimensional embedding is optimized by minimizing the cross-entropy between the high-dimensional and low-dimensional representations, ensuring the structural similarities are preserved. UMAP \cite{umap_reference} uses stochastic gradient descent to iteratively refine the positions of the points, effectively embedding the high-dimensional data in a low-dimensional space while maintaining global structures, and these authors also claimed that it does not guarantee the local structure.


UMAP \cite{umap_reference} outperforms t-SNE \cite{tsne} in preserving global structures and is scalable for large datasets. This approach provides good visualizations and is computationally less expensive than t-SNE, making it suitable for real-time datasets.


\section{Exploring Laplacian Eigenmaps and LAPtSNE}

Laplacian Eigenmaps \cite{belkin2001laplacian}, a spectral graph-based technique, preserves the local structure of the data. It achieves this by creating a graph with nodes representing data points and edges reflecting their neighbourhood associations. The approach then computes the Laplacian matrix for the created graph and, based on these values, further calculates eigenvalues to determine a lower-dimensional data embedding. This embedding minimizes the cost function by ensuring that the closeness between data points in the high-dimensional space remains the same in the low-dimensional space. Laplacian Eigenmaps are highly effective in capturing the multidimensional structure of data and particularly helpful for clustering the data in visualization space. 

Sun et al., \cite{sun2023laplacian} developed the LAPtSNE approach that uses the base t-SNE algorithm. This approach preserves both local and global structures in high-dimensional space, resulting in compact cluster visualization in low-dimensional embedding. The process begins by applying t-SNE to compute pairwise affinities in the high-dimensional data. Further, the graph laplacian matrix is calculated based on the affinity matrix of t-SNE. The \textit{`d'} number of the smallest eigenvectors of the Laplacian matrix is used to initialize the low-dimensional embedding. 

In each iteration of LAPtSNE, low-dimensional probabilities are computed using a student-t distribution, similar to the standard t-SNE process. The LAPtSNE algorithm operates in an iterative manner, continually refining the low-dimensional embedding. The gradient of the KL divergence between high and low-dimensional affinities is estimated and regularized. The low-dimensional embeddings are updated with the calculated values to minimize the model’s KL divergence and the regularization factor to preserve neighbourhood relations. This iterative process is crucial in refining the low-dimensional embedding while preserving the global structure and cluster integrity.


An approach like t-SNE \cite{tsne} preserves the local structure but claims to retain the global structure partially to some extent. UMAP \cite{umap_reference} focuses more on preserving the global structure but does not guarantee the local structure. LAPtSNE \cite{sun2023laplacian} can preserve both local and global structures with cluster visualization quality, but it is computationally intensive for high-dimensional datasets.  Therefore, there is a need to discover a computationally efficient NLDR approach that can preserve both local and global structures and produce high-quality cluster visualizations for high-dimensional data.
